\section{Problem Formulation}

Our goal is to determine whether a shared structural inductive bias can be
recovered from a collection of successfully trained models and transferred
to unseen tasks. We formalize the construction of the generative-model
training set, the resulting structural prior, and the criterion used to
evaluate its transferability.

\paragraph*{Meta-learning environment}

Let $\mathcal{T}$ be a space of tasks. Each task $\tau\in\mathcal{T}$ defines
a data distribution $P_\tau$ and a loss function $\mathcal{L}_\tau$.
Meta-training and meta-test tasks are sampled from distributions
$\mu_{\mathrm{tr}}$ and $\mu_{\mathrm{test}}$, respectively. Meta-test tasks
are held out from all stages of generative-model training.

For each task $\tau$, we denote its training, validation, and test datasets
by $\mathcal{D}^{\mathrm{tr}}_\tau$,
$\mathcal{D}^{\mathrm{val}}_\tau$, and
$\mathcal{D}^{\mathrm{test}}_\tau$, respectively. These datasets consist of
independent samples from $P_\tau$.

A model for task $\tau$ is denoted by $f_{a,w}$, where
$a\in\mathcal{A}$ represents its discrete or continuous structure and
$w\in\mathcal{W}_a$ denotes its learnable parameters. Candidate structures
are initially sampled from a proposal distribution $p_0$.

Given a structure $a$, its parameters are obtained by minimizing the
training loss:
$$
\widehat{w}_\tau(a)
\in
\operatorname*{arg\,min}_{w\in\mathcal{W}_a}
\mathcal{L}_\tau
\left(
f_{a,w};
\mathcal{D}^{\mathrm{tr}}_\tau
\right).
$$
The resulting solution to task $\tau$ is
$$
s_\tau(a)
=
\left(
a,\widehat{w}_\tau(a)
\right).
$$

\paragraph*{Successful solution representations}

Let
$\tau_1,\ldots,\tau_T\sim\mu_{\mathrm{tr}}$
be the meta-training tasks. For each task $\tau_j$, we independently sample
candidate structures
$a_{j,1},\ldots,a_{j,N}\sim p_0$.

The validation loss of a solution is defined as
$$
V_\tau(a)
=
\mathcal{L}_\tau
\left(
f_{a,\widehat{w}_\tau(a)};
\mathcal{D}^{\mathrm{val}}_\tau
\right).
$$
A solution is considered successful if
$V_\tau(a)\leq\varepsilon_\tau$, where $\varepsilon_\tau$ is a
task-specific threshold. Thus, the successful solutions for task $\tau_j$
form the set
$$
\mathcal{S}^{\mathrm{succ}}_{\tau_j}
=
\left\{
s_{\tau_j}(a_{j,i})
\;\middle|\;
V_{\tau_j}(a_{j,i})\leq\varepsilon_{\tau_j},
\quad i=1,\ldots,N
\right\}.
$$

A representation map $\psi$ assigns each solution $s_\tau(a)$ a
representation $\psi(s_\tau(a))\in\mathcal{X}$. This representation may
depend on the structure $a$, the trained parameters
$\widehat{w}_\tau(a)$, or both. The training dataset for the generative model
is then
$$
\mathcal{D}_{\mathrm{gen}}
=
\left\{
\psi(s)
\;\middle|\;
s\in\mathcal{S}^{\mathrm{succ}}_{\tau_j},
\quad j=1,\ldots,T
\right\}.
$$

\paragraph*{Learned structural prior}

Let $p_z$ be a distribution on a latent space $\mathcal{Z}$, for example,
$p_z=\mathcal{N}(\mathbf{0},\mathbf{I})$. A generative model
$G_\theta:\mathcal{Z}\rightarrow\mathcal{X}$ is trained on
$\mathcal{D}_{\mathrm{gen}}$. A mapping
$\pi:\mathcal{X}\rightarrow\mathcal{A}$ transforms a generated
representation into a model structure:
$$
\widetilde{a}
=
\pi\left(G_\theta(z)\right),
\qquad
z\sim p_z.
$$
The distribution induced over model structures is
$$
q_\theta
=
\left(\pi\circ G_\theta\right)_{\#}p_z,
$$
where $(\cdot)_{\#}$ denotes the pushforward of a distribution. We refer to
$q_\theta$ as the learned structural prior and hypothesize that it captures
structural inductive biases shared across the meta-training tasks.

\paragraph*{Out-of-distribution evaluation}

For a task $\tau$ and a distribution $q$ over model structures, we define
the expected test risk as
$$
\mathcal{R}_\tau(q)
=
\mathbb{E}_{a\sim q}
\left[
\mathcal{L}_\tau
\left(
f_{a,\widehat{w}_\tau(a)};
\mathcal{D}^{\mathrm{test}}_\tau
\right)
\right].
$$
The meta-test transfer gain of the learned prior relative to the initial
proposal distribution is defined as
$$
\Delta_{\mu_{\mathrm{test}}}(q_\theta,p_0)
=
\mathbb{E}_{\tau\sim\mu_{\mathrm{test}}}
\left[
\mathcal{R}_\tau(p_0)
-
\mathcal{R}_\tau(q_\theta)
\right].
$$
A positive transfer gain indicates that structures sampled from
$q_\theta$ outperform those sampled from $p_0$ on unseen tasks. Since the
parameters $\widehat{w}_\tau(a)$ are trained separately for each sampled
structure, this improvement is attributable to the learned structural prior
rather than to transferred model parameters.